\section{Introduction}
%\note{
%Answer the following questions
%\begin{itemize}
%  \item What is the problem you are solving?
%  \begin{itemize}
%    \item Why is it important?
%    \item How prevalent is it?
%  \end{itemize}
%  \item How is it solved currently, and what are the limitations of current methods?
%  \item What is your approach? How does it overcome the limitations above?
%  \item How do you know your approach works? (Evaluation)
%  \item What difference does it make?
%\end{itemize}
%
%}

The dramatic growth of behavioral data has created new opportunities for insights into individual human and social behavior through computational methods. However, computational analysis and modeling of large-scale data also presents many challenges. Real-world behavioral data is extremely noisy. To uncover hidden patterns and build predictive models, scientists aggregate data over the entire population, which tends to improve statistics and signal-to-noise ratio. Real-world data is also highly heterogeneous, i.e., composed of subgroups that vary widely in size and behavior. This heterogeneity can severely bias analysis. %estimates of causal effects.
One example of such bias is Simpson's paradox. According to the paradox, which is often encountered in economics, demographics and clinical research, an association observed in data that has been aggregated over a population may be quite different from, and even opposite to, those of the underlying subgroups. Failure to take this effect into account can distort conclusions drawn from data.
One of the better known examples of Simpson's paradox comes from a study of gender bias in graduate admissions~\cite{Bickel1975}. In the aggregate admissions data there appears to be a statistically significant bias against women: a smaller fraction of female applicants is admitted for graduate studies. However, when admissions data is disaggregated by department, women have parity and even a slight advantage over men in some departments. The paradox arises because departments preferred by female applicants have lower admissions rates for both genders.


Simpson's paradox also exists in the analysis of trends. When measuring an outcome changes as a function of an independent variable, the characteristics of the population over which the trend is measured may change as a function of the independent variable due to survivor bias. As a result, the data may appear to exhibit a trend, which disappears or reverses when the data is disaggregated. %The trends discovered in disaggregated are more robust and more likely to describe individual behavior than the trends found in aggregate data.
For example, it may appear that repeated exposures to information or hashtags on social media make an individual less likely to share it with his or her followers~\cite{Romero11www}. In fact, the opposite is true: the more people are exposed to information, the more likely they are to share it with followers~\cite{Lerman2016futureinternet}. However, those who follow many others, and are likely to be exposed to a meme or a hashtag multiple times, are less responsive overall, due to high volume of other information they receive~\cite{Hodas12socialcom}. Their reduced susceptibility to exposures biases the aggregate response, leading to wrong conclusions about behavior. Once disaggregated based on the volume of information received, a clearer pattern of exposure response emerges, that is more predictive of the actual response~\cite{Hodas14srep}.
Despite accumulating evidence that Simpson's paradox %could distort the associations found in
affects inference of trends in social and behavioral data~\cite{Singer2016plosone,Barbosa2016,KootiA2017www,Agarwal2017quitting}, researchers do not routinely test for it in their studies.

% What we are doing
We describe a method for detecting Simpson's paradox in the analysis of trends in heterogeneous data. Our method identifies trends in the aggregate data that disappear or reverse when the data is disaggregated into its underlying subgroups.
%Our method automatically identifies covariates in data that could be used to disaggregate it into more homogeneous subgroups.
Our statistical approach identifies pairs of variables---what we call Simpson's pairs---such that a trend in some outcome as a function of the first variable that is observed in the aggregate data  disappears or reverses when the data is disaggregated on the second variable.

% Evaluation
We use the proposed approach to study data collected from Stack Exchange, a popular question-answering platform. Specifically, we analyze factors affecting answerer performance, measured as the likelihood the answer provided by the answerer will be accepted by the asker as the best answer to his or her question. We construct a variety of variables describing features of answers and answerers, and study how the outcome---in our case answer acceptance---depends on these variables. We compare results of regression analysis on aggregated and disaggregated data to identify Simpson's pairs. Our analysis identifies several cases of Simpson's paradox in Stack Exchange data. In addition to a known effect that describes deterioration of answerer performance over the course of a session, our method identifies several new cases of Simpson's paradox. These paradoxes yield new insights into answerer performance on Stack Exchange. For example, creating a new feature to describe answerers, which combines variables of a Simpson's pair, leads to a more robust proxy of answerer performance.

% Summary
Social data often comes from populations comprised of subgroups that differ in their behavior. Presence of a Simpson's paradox in such data can indicate interesting patterns~\cite{fabris2000discovering}, including important behavioral differences. The paradox suggests that subgroups within the population under study systematically differ in their behavior, and these differences are large enough to affect aggregate trends. In such cases the trends discovered in disaggregated data are more likely to describe---and predict---individual behavior than the trends found in aggregated data. Thus, to build more robust models of behavior, data scientists need to identify confounding variables which could affect their conclusions. The method presented in this paper provides a simple framework for identifying such variables by systematically searching for Simpson's trend paradox in data.
Simpson's paradox as a tool to discover new important features of data. e.g., we discover a better measure of answerer quality.


The rest of the paper is organized as follows. First, we present background, including multiple observations of Simpson's paradox in a variety of disciplies (Sec.~\ref{sec:related}). Then we describe our methodology for detecting Simpson's paradox by identifying covariates in data, and analyze the paradox mathematically to gain more insight into its origins (Sec.~\ref{sec:methods}). Finally, we apply our method to real-world data from the question-answering site Stack Exchange, and demonstrate its ability to automatically identify novel cases of Simpson's paradox (Sec.~\ref{sec:results}). We conclude with the discussion of the implications of our results.

